\section{The two atomic edges and the logarithmic corner}

\begin{theorem}[Boundary-layer scale and weak convergence]
\label{subunit:thm:boundarylayer}
Fix $0<b<1$, let $a=1-b+\varepsilon$ with $\varepsilon>0$, and
let $T_\varepsilon=-\log r_{a,b}$. Define
\[
 \tau_\varepsilon=
 \left(\frac{\Gamma(1-b+\varepsilon)}
 {(1-b)(-\zeta(b))\Gamma(\varepsilon)}\right)^{1/(1-b)}.
\]
Then, as $\varepsilon\downarrow0$,
\begin{equation}\label{subunit:eq:layerexact}
 T_\varepsilon=\tau_\varepsilon
 \left(1+O_b\left(\frac{\tau_\varepsilon}{\varepsilon}\right)\right),
 \qquad 0<T_\varepsilon<\tau_\varepsilon,
\end{equation}
and in particular
\begin{equation}\label{subunit:eq:layerleading}
 1-r_{1-b+\varepsilon,b}\sim
 \left(\frac{\Gamma(1-b)}{(1-b)(-\zeta(b))}
       \varepsilon\right)^{1/(1-b)}.
\end{equation}
The measures converge weakly to the critical measure in
Theorem~\ref{subunit:thm:domain}, and their positive masses converge
to $(1-b)^{-1}$ while their positive supports shrink to $u=1$.
\end{theorem}

\begin{proof}
In \eqref{subunit:eq:rescaledLowB}, put
$A_\varepsilon=((1-b)\Gamma(\varepsilon))^{-1}$ and
$B_\varepsilon=(-\zeta(b))/\Gamma(1-b+\varepsilon)$.
The zero equation is
\[
 A_\varepsilon-B_\varepsilon T^{1-b}-T C_\varepsilon(T)=0,
 \quad
 C_\varepsilon(T)=\frac1{\Gamma(a)\Gamma(b)}
 \int_0^1(1-v)^{a-1}v^{b-1}q(Tv)\,dv.
\]
Since $1/2<q<1$ and $a+b=1+\varepsilon$,
\[
 \frac1{2\Gamma(1+\varepsilon)}
 <C_\varepsilon(T)<\frac1{\Gamma(1+\varepsilon)}.
\]
By definition $B_\varepsilon\tau_\varepsilon^{1-b}
=A_\varepsilon$, so the equation at $T=\tau_\varepsilon$ has
strictly negative left side and $T_\varepsilon<\tau_\varepsilon$.
Furthermore
\[
 0\leq1-(T_\varepsilon/\tau_\varepsilon)^{1-b}
 =\frac{T_\varepsilon C_\varepsilon(T_\varepsilon)}
       {A_\varepsilon}
 \leq\frac{(1-b)\tau_\varepsilon}{\varepsilon}.
\]
The final equality in the upper bound uses
$\Gamma(1+\varepsilon)=\varepsilon\Gamma(\varepsilon)$.
Now $\tau_\varepsilon=O_b(\varepsilon^{1/(1-b)})$, hence
$\tau_\varepsilon/\varepsilon\to0$. The last inequality proves
\eqref{subunit:eq:layerexact}. Expanding the Gamma ratio and using
$1-e^{-T}\sim T$ proves \eqref{subunit:eq:layerleading}.

For weak convergence, the positive term in
\eqref{subunit:eq:generalK}, regarded as a measure in the $T$
coordinate, is
\[
 \frac1{1-b}\frac{e^{-T}T^{\varepsilon-1}}
 {\Gamma(\varepsilon)}\,dT.
\]
It is $(1-b)^{-1}$ times a Gamma$(\varepsilon,1)$ law. Its mean
is $\varepsilon$, so it converges weakly to mass $(1-b)^{-1}$
at $T=0$, equivalently at $u=1$. The two negative components in
\eqref{subunit:eq:generalK} converge in $L^1(e^{-T}\,dT)$ to
those in \eqref{subunit:eq:criticaldensity}. For the power term,
dominated convergence holds with an exponent bounded above $-1$.
For the convolution, either use the rescaled beta integral and the
same domination, or use $L^1$ continuity of the Gamma densities and
the integrable convolution factor $e^{-t}t^{b-1}q(t)$.

The unique positive support is $0<T<T_\varepsilon$, which shrinks
to zero. The Gamma$(\varepsilon,1)$ mass in this interval tends to
one: since $T_\varepsilon$ is a fixed power of $\varepsilon$ to
leading order,
\[
 \frac1{\Gamma(\varepsilon)}\int_0^{T_\varepsilon}
 e^{-T}T^{\varepsilon-1}\,dT
 =\frac{T_\varepsilon^\varepsilon}{\Gamma(1+\varepsilon)}(1+o(1))
 \longrightarrow1.
\]
The negative components have vanishing mass on the shrinking interval
by their $L^1$ convergence to an integrable limit. Thus the actual
positive mass of the signed density tends to $(1-b)^{-1}$, completing
the proof.
\end{proof}

\subsection{The other edge and the logarithmic corner}

The second atomic edge has a complementary power law. The corner $b=1$
has an exponential scale and no finite limiting measure.

\begin{proposition}[Two further endpoint scales]
\label{subunit:prop:otheredges}
For fixed $b>1$, as $a\downarrow0$,
\begin{equation}\label{subunit:eq:otheredge}
 1-r_{a,b}\sim\bigl(\Gamma(b)\zeta(b)a\bigr)^{1/(b-1)}.
\end{equation}
The measures converge weakly to $\nu_{0,b}$, and their positive masses
converge to $\zeta(b)$. For $b=1$, set
$\sigma_a=\exp(\gamma+\psi(a))$. Then
\begin{equation}\label{subunit:eq:logcorner}
 -\log r_{a,1}=\sigma_a\bigl(1+O(\sigma_a/a)\bigr),
 \qquad 0<-\log r_{a,1}<\sigma_a,
\end{equation}
where
\[
 \sigma_a=\exp\bigl(-1/a+\zeta(2)a-\zeta(3)a^2+O(a^3)\bigr).
\]
If $M_{a,1}$ denotes either signed mass of $\nu_{a,1}$, then
\begin{equation}\label{subunit:eq:divergingmass}
 M_{a,1}\sim\frac1{e a}.
\end{equation}
Thus the finite signed measures have unbounded total variation at this
corner, consistently with the impossibility of a finite measure at
$a=0,b=1$.
\end{proposition}

\begin{proof}
For $b>1$, the equation obtained from
\eqref{subunit:eq:generalK} after division by $T^{a-1}$ is
\[
 0=\frac{\zeta(b)}{\Gamma(a)}
 -\frac{T^{b-1}}{(b-1)\Gamma(a+b-1)}
 -\frac{T^b}{\Gamma(a)\Gamma(b)}
       \int_0^1(1-v)^{a-1}v^{b-1}q(Tv)\,dv.
\]
The last integral, after its Gamma normalization, lies between
$1/(2\Gamma(a+b))$ and $1/\Gamma(a+b)$. The ratio of this
last term to the preceding negative term is therefore $O_b(T)$,
uniformly for small positive $a$. The balancing point of the first
two terms is
\[
 \left(\frac{(b-1)\Gamma(a+b-1)\zeta(b)}{\Gamma(a)}\right)^{1/(b-1)}
 \sim\bigl(\Gamma(b)\zeta(b)a\bigr)^{1/(b-1)}.
\]
The strict monotonicity of the rescaled density, or the same two-sided
comparison used in Theorem~\ref{subunit:thm:boundarylayer}, makes the
actual zero asymptotic to this point. The positive term is $\zeta(b)$
times a Gamma$(a,1)$ law in the $T$ coordinate. The negative term in
\eqref{subunit:eq:largeB}, after multiplication by $e^{-T}$, is the
convolution of that Gamma probability density with
$e^{-t}t^{b-1}/(\Gamma(b)(1-e^{-t}))$, an integrable function of mass
$\zeta(b)$. As $a\downarrow0$, this convolution converges in $L^1$
to the latter function; Gamma$(a,1)$ is an approximate identity since
its mean tends to zero. This proves weak convergence. As before,
the Gamma mass below the crossing tends to one, since $a\log T\to0$,
and the limiting negative density is integrable. Hence the positive
mass tends to $\zeta(b)$.

For $b=1$, put
$J_a(T)=\int_0^1(1-v)^{a-1}q(Tv)\,dv$, so that
$1/(2a)<J_a(T)<1/a$. At the crossing $T_a$, formula
\eqref{subunit:eq:bOne} says
\[
 \log(\sigma_a/T_a)=T_aJ_a(T_a),\qquad
 0<\log(\sigma_a/T_a)<\sigma_a/a.
\]
This gives \eqref{subunit:eq:logcorner}. The digamma expansion
$\psi(a)=-1/a-\gamma+\zeta(2)a-\zeta(3)a^2+O(a^3)$ gives
the stated scale.

Finally the positive mass is the integral of
$e^{-T}T^{a-1}[\log(\sigma_a/T)-T J_a(T)]/\Gamma(a)$
over $0<T<T_a$. Omitting the factor $e^{-T}$, the $TJ_a$ term,
and the vanishing interval between $T_a$ and $\sigma_a$ changes
the result by a relative $o(1)$. To see this explicitly, the first
omission has relative error at most $\sigma_a$; the $TJ_a$ integral
is at most $\sigma_a^{a+1}/(a(a+1)\Gamma(a))$, whereas the main
integral is $\sigma_a^a/(a^2\Gamma(a))$. For the removed interval,
write $T=\sigma_a v$ and use
$1-T_a/\sigma_a=O(\sigma_a/a)$; its relative contribution is
$O(\sigma_a^2)$. Therefore
\[
 M_{a,1}\sim\frac{\sigma_a^a}{a^2\Gamma(a)}
 \sim\frac1{ea},
\]
as asserted.
\end{proof}
